%% A realistic Springer LNCS paper. Springer ships no sample .tex on CTAN, so
%% this follows the shape llncs.cls itself defines. Used to test venue-shift.
\documentclass[runningheads]{llncs}

\usepackage{graphicx}
\usepackage{booktabs}

\begin{document}

\title{Trust Calibration in Highly Automated Driving\thanks{Supported by the
German Research Foundation.}}
\titlerunning{Trust Calibration in Automated Driving}

\author{Mark Colley\inst{1}\orcidID{0000-0001-5207-5029} \and
Jane Doe\inst{2} \and
Alex Roe\inst{1}}

\authorrunning{M. Colley et al.}

\institute{Institute of Media Informatics, Ulm University, Ulm, Germany\\
\email{mark.colley@uni-ulm.de} \and
Department of Computer Science, Example University, Boston, USA\\
\email{jane.doe@example.edu}}

\maketitle

\begin{abstract}
Automated vehicles must communicate their intent to build appropriate trust.
We report a within-subjects study with 24 participants examining how takeover
requests shape trust over repeated exposure.

\keywords{automated driving \and trust \and takeover request \and user study}
\end{abstract}

\section{Introduction}
Trust in automated vehicles is neither uniformly good nor uniformly
bad~\cite{lee2004trust}. We examine how it develops across repeated takeovers.

\begin{figure}[t]
  \centering
  \includegraphics[width=\textwidth]{figures/apparatus}
  \caption{The driving simulator used in the study.}
  \label{fig:apparatus}
\end{figure}

\section{Method}
We ran a within-subjects experiment with 24 participants.

\begin{table}[t]
  \caption{Participant demographics.}
  \label{tab:demographics}
  \begin{tabular}{lr}
    \toprule
    Measure & Value \\
    \midrule
    Participants & 24 \\
    Mean age & 27.4 \\
    \bottomrule
  \end{tabular}
\end{table}

\section{Results}
Trust increased over exposure, $F(2,46) = 8.1$, $p < .001$.

\section{Conclusion}
Repeated exposure calibrates trust.

\bibliographystyle{splncs04}
\bibliography{references}

\end{document}
